\clearpage
\renewcommand{\sectioncolor}{bio}
\section{The route, through books you already own}
\markboth{The route}{}

\begin{center}
\begin{tikzpicture}[font=\footnotesize,
  st/.style={rounded corners=4pt,draw=#1,line width=1.2pt,fill=#1!8,inner sep=8pt,
             text width=3.55cm,align=left,minimum height=4.4cm},
  ar/.style={-{Stealth[length=5.5pt,width=4.5pt]},line width=1.2pt,draw=bio!75}]
 \node[st=bio]  (s1) at (0,0)     {\textbf{\color{bio}STAGE 1}\\\textbf{Operators \& the trace}\\[4pt]
   {\scriptsize Axler §5 eigenvalues\\ §6 inner products\\ \textbf{§7.B spectral theorem}\\
    \textbf{§7.C positive operators}\\ \textbf{§10.A trace}}\\[4pt]
   {\scriptsize + Hoffman \& Kunze\\ + Olver ch.\ 8--10}};
 \node[st=nano] (s2) at (4.25,0)  {\textbf{\color{nano}STAGE 2}\\\textbf{Integration $\to$ measure}\\[4pt]
   {\scriptsize Apostol Vol.~2 ch.~11\\ Spivak, \textit{Calc.\ on Manifolds}\\[3pt]
    then \textbf{Kuttler ch.~20--24}: measures, Lebesgue integral, $L^p$, Riesz}};
 \node[st=sand] (s3) at (8.50,0)  {\textbf{\color{sand!80!black}STAGE 3}\\\textbf{Saddle point}\\[4pt]
   {\scriptsize Radozycki Part III:\\ contour integrals,\\ Laurent series,\\ \textbf{the residue theorem},\\
    Fourier series}\\[4pt]
   {\scriptsize this is how $Z$ is \emph{evaluated} asymptotically}};
 \node[st=quantum] (s4) at (12.75,0) {\textbf{\color{quantum}STAGE 4}\\\textbf{Probability}\\[4pt]
   {\scriptsize Apostol Vol.~2\\ ch.~13 set functions\\ ch.~14 probability}\\[4pt]
   {\scriptsize \textcolor{alert}{\textbf{the shelf runs thin here}} --- see §4}};
 \foreach \a/\b in {s1/s2,s2/s3,s3/s4}{\draw[ar] (\a)--(\b);}
\end{tikzpicture}
\end{center}
\vspace{2pt}
\begin{center}\begin{minipage}{0.93\textwidth}\footnotesize\color{slate}
\textbf{Figure 3.}\; Four stages, all but the last fully covered by books already in the tree.
\end{minipage}\end{center}

\subsection{Stage 1 in detail --- the four pages that matter}

I checked the actual contents of your copy of Axler. The three sections you need are:

\begin{center}\small
\begin{tabular}{@{}>{\raggedright\arraybackslash}p{4.4cm}c>{\raggedright\arraybackslash}p{9.6cm}@{}}
\toprule
\rowcolor{bio!13}
\textbf{\color{bio}Axler section} & \textbf{p.} & \textbf{\color{bio}Why $Z=\operatorname{Tr}(e^{-\beta\hat H})$ needs it}\\
\midrule
\textbf{7.B The Spectral Theorem}\newline{\scriptsize Complex / Real} & 217 / 219 &
 This is what lets you \emph{define} $e^{-\beta\hat H}$ at all --- you exponentiate the eigenvalues,
 which requires an orthonormal eigenbasis to exist.\\
\textbf{7.C Positive Operators} & 225 &
 $e^{-\beta\hat H}$ \textbf{is} a positive operator. This is why $Z>0$, and why
 $e^{-\beta U}/Z$ is a probability density rather than merely a function.\\
\textbf{10.A Trace} & 296--299 &
 ``Trace: a connection between operators and matrices'' --- basis independence is what makes
 $\operatorname{Tr}(e^{-\beta\hat H})=\sum_n e^{-\beta E_n}$ a definition rather than a coincidence.\\
\bottomrule
\end{tabular}
\end{center}

\begin{haveit}[title={You can do Stage 1 this week}]
Read Axler §7.B, §7.C and §10.A --- about 30 pages with the exercises. Then go straight to §5 of this
dossier and compute a real partition function with nothing but those three results. That is the whole
argument for starting at Face 3.
\end{haveit}
